\documentclass[11pt,a4paper]{article}
\usepackage[utf8]{inputenc}
\usepackage[T1]{fontenc}
\usepackage[ngerman]{babel}
\usepackage{geometry}
\usepackage{listings}
\usepackage{xcolor}
\usepackage{hyperref}
\usepackage{booktabs}
\usepackage{fancyhdr}

\geometry{margin=2.5cm}

% Code-Stil
\definecolor{codebg}{RGB}{248,248,248}
\definecolor{codegreen}{RGB}{40,160,40}
\definecolor{codegray}{RGB}{128,128,128}
\definecolor{codeblue}{RGB}{40,40,200}

\lstdefinestyle{code}{
    backgroundcolor=\color{codebg},
    basicstyle=\ttfamily\small,
    keywordstyle=\color{codeblue}\bfseries,
    stringstyle=\color{codegreen},
    commentstyle=\color{codegray}\itshape,
    breaklines=true,
    frame=single,
    framerule=0.5pt,
    numbers=left,
    numberstyle=\tiny\color{codegray},
    tabsize=4
}

\pagestyle{fancy}
\fancyhf{}
\fancyhead[L]{\textsc{Local Supermemory}}
\fancyhead[R]{\thepage}

\title{\textbf{Local Supermemory}\\[0.5em]\large Entwicklerhandbuch v1.0}
\author{granaria}
\date{Februar 2026}

\begin{document}
\maketitle
\tableofcontents
\newpage

%──────────────────────────────────────────────────────────────
\section{Übersicht}
%──────────────────────────────────────────────────────────────

Local Supermemory ist ein MCP-Server, der Claude Desktop ein persistentes Gedächtnis gibt. Alle Daten bleiben lokal.

\subsection{Architektur}

\begin{verbatim}
┌─────────────────┐     ┌──────────────┐
│  Claude Desktop │────▶│  MCP Server  │  (stdio)
└─────────────────┘     └──────┬───────┘
                               │
              ┌────────────────┼────────────────┐
              │                │                │
        ┌─────▼─────┐    ┌─────▼─────┐    ┌─────▼─────┐
        │  ChromaDB │    │  SQLite   │    │  Ollama   │
        │ (Vektoren)│    │(Metadaten)│    │ (Profil)  │
        └───────────┘    └───────────┘    └───────────┘
\end{verbatim}

\textbf{Komponenten:}
\begin{itemize}
    \item \textbf{ChromaDB} -- Vektordatenbank für semantische Ähnlichkeitssuche
    \item \textbf{SQLite} -- Relationale Daten: Memories, Projekte, Profil-Cache
    \item \textbf{Ollama} -- Lokales LLM für Profil-Generierung (optional)
\end{itemize}

\subsection{Datenspeicherung}

\begin{lstlisting}[style=code,language=bash]
~/.local-supermemory/
+-- chroma/           # ChromaDB Vektoren
|   +-- chroma.sqlite3
+-- memories.db       # SQLite Metadaten
\end{lstlisting}

%──────────────────────────────────────────────────────────────
\section{Installation}
%──────────────────────────────────────────────────────────────

\begin{lstlisting}[style=code,language=bash]
git clone https://github.com/granaria/local-supermemory.git
cd local-supermemory
pip install -e .
\end{lstlisting}

\textbf{Abhängigkeiten:} \texttt{mcp}, \texttt{chromadb}, \texttt{httpx}

\subsection{Claude Desktop Config}

Datei: \texttt{\textasciitilde/Library/Application Support/Claude/claude\_desktop\_config.json}

\begin{lstlisting}[style=code]
{
  "mcpServers": {
    "local-supermemory": {
      "command": "python3",
      "args": ["-m", "local_supermemory.server"]
    }
  }
}
\end{lstlisting}


%──────────────────────────────────────────────────────────────
\section{MCP Tools}
%──────────────────────────────────────────────────────────────

\subsection{memory}

Speichert oder löscht eine Memory.

\begin{tabular}{llll}
\toprule
\textbf{Parameter} & \textbf{Typ} & \textbf{Default} & \textbf{Beschreibung} \\
\midrule
\texttt{action} & string & \texttt{"save"} & \texttt{"save"} oder \texttt{"forget"} \\
\texttt{content} & string & -- & Memory-Text (max 200KB) \\
\texttt{project} & string & \texttt{"default"} & Projekt-Scope \\
\bottomrule
\end{tabular}

\vspace{1em}
\textbf{Interner Ablauf:}
\begin{enumerate}
    \item Content-Hash (SHA256[:16]) als ID generieren
    \item ChromaDB: Dokument mit Embedding speichern
    \item SQLite: Metadaten persistieren
    \item Profil-Cache invalidieren
\end{enumerate}

\begin{lstlisting}[style=code,language=Python,caption={store.py -- save()}]
def save(self, content: str, project: str = "default"):
    mid = self._gen_id(content)  # SHA256[:16]
    coll = self._get_collection(project)
    
    # Upsert in ChromaDB
    existing = coll.get(ids=[mid])
    if existing['ids']:
        coll.update(ids=[mid], documents=[content])
    else:
        coll.add(ids=[mid], documents=[content])
    
    # SQLite: INSERT ... ON CONFLICT DO UPDATE
    # Cache invalidieren
    conn.execute("DELETE FROM profile_cache WHERE project=?")
\end{lstlisting}


\subsection{recall}

Semantische Suche mit optionalem Profil.

\begin{tabular}{llll}
\toprule
\textbf{Parameter} & \textbf{Typ} & \textbf{Default} & \textbf{Beschreibung} \\
\midrule
\texttt{query} & string & -- & Suchanfrage \\
\texttt{project} & string & \texttt{"default"} & Projekt-Scope \\
\texttt{include\_profile} & bool & \texttt{true} & Profil anhängen \\
\texttt{n\_results} & int & 15 & Max. Ergebnisse \\
\bottomrule
\end{tabular}

\vspace{1em}
\textbf{Ablauf:}
\begin{enumerate}
    \item Query embedden (MiniLM-L6-v2, 384d Vektoren)
    \item Cosine-Similarity-Suche in ChromaDB Collection
    \item Optional: Profil aus Cache oder via Ollama generieren
    \item Ergebnisse mit Similarity-Score (0--100\%) zurückgeben
\end{enumerate}

\begin{lstlisting}[style=code,language=Python,caption={store.py -- recall()}]
def recall(self, query: str, project: str = "default", n: int = 15):
    coll = self._get_collection(project)
    res = coll.query(query_texts=[query], n_results=n)
    
    memories = []
    for i, doc in enumerate(res['documents'][0]):
        dist = res['distances'][0][i]  # Cosine: [0,2]
        memories.append({
            "content": doc,
            "similarity": round((1 - dist) * 100),  # -> %
            "id": res['ids'][0][i]
        })
    return memories
\end{lstlisting}

\subsection{Weitere Tools}

\begin{description}
    \item[list\_projects] Listet alle Projekte mit Memory-Count
    \item[stats] Statistiken: Gesamtzahl, Projekte, Pfad
    \item[whoami] Generiert Profil aus letzten 10 Memories via Ollama
\end{description}


%──────────────────────────────────────────────────────────────
\section{Datenbank-Schema}
%──────────────────────────────────────────────────────────────

\begin{lstlisting}[style=code,language=SQL,caption={SQLite Schema}]
CREATE TABLE memories (
    id TEXT PRIMARY KEY,        -- SHA256[:16] vom Content
    content TEXT NOT NULL,
    project TEXT DEFAULT 'default',
    created_at TEXT NOT NULL,   -- ISO 8601
    updated_at TEXT NOT NULL,
    metadata TEXT               -- JSON
);

CREATE TABLE projects (
    name TEXT PRIMARY KEY,
    description TEXT,
    created_at TEXT NOT NULL
);

CREATE TABLE profile_cache (
    project TEXT PRIMARY KEY,
    profile TEXT NOT NULL,      -- LLM-generierter Text
    generated_at TEXT NOT NULL
);

CREATE INDEX idx_memories_project ON memories(project);
\end{lstlisting}

\textbf{Design-Entscheidungen:}
\begin{itemize}
    \item \textbf{Content-Hash als ID} -- Deduplizierung: gleicher Text = gleiche ID
    \item \textbf{Profil-Cache} -- Vermeidet LLM-Aufrufe, invalidiert bei Änderung
    \item \textbf{Dual Storage} -- ChromaDB für Vektoren, SQLite für Queries
\end{itemize}


%──────────────────────────────────────────────────────────────
\section{Profil-Engine}
%──────────────────────────────────────────────────────────────

Aggregiert Memories zu strukturiertem Benutzerprofil via Ollama.

\begin{lstlisting}[style=code,language=Python,caption={profile.py}]
PROMPT = """Analysiere folgende Benutzer-Memories...
Erstelle Profil mit: Identitaet, Hardware/Software, 
Projekte, Praeferenzen, Kontakte.
Nur vorhandene Infos, keine Spekulationen."""

class ProfileEngine:
    def __init__(self, url="http://localhost:11434", 
                 model="qwen2.5:32b"):
        self.url = url
        self.model = model
    
    async def generate(self, memories: list) -> str:
        text = "\n".join([f"- {m['content']}" for m in memories[:50]])
        r = await client.post(f"{self.url}/api/generate", json={
            "model": self.model,
            "prompt": PROMPT.format(memories=text),
            "options": {"temperature": 0.3}
        })
        return r.json()["response"]
\end{lstlisting}

\textbf{Config:} Model änderbar in \texttt{profile.py}. Fallback ohne Ollama: einfache Liste.

%──────────────────────────────────────────────────────────────
\section{ChromaDB Details}
%──────────────────────────────────────────────────────────────

\begin{lstlisting}[style=code,language=Python,caption={ChromaDB Init}]
self.chroma = chromadb.PersistentClient(
    path=str(self.data_dir / "chroma"),
    settings=Settings(anonymized_telemetry=False)
)

def _get_collection(self, project: str = "default"):
    name = f"memories_{project.replace('-', '_')}"
    return self.chroma.get_or_create_collection(
        name=name, 
        metadata={"hnsw:space": "cosine"}
    )
\end{lstlisting}

\textbf{Embedding:} \texttt{all-MiniLM-L6-v2} (384d), automatisch bei \texttt{add()}

\textbf{HNSW-Index:} Hierarchical Navigable Small World -- ANN-Suche, Cosine Distance $\in [0,2]$


%──────────────────────────────────────────────────────────────
\section{MCP Server}
%──────────────────────────────────────────────────────────────

\begin{lstlisting}[style=code,language=Python,caption={server.py Struktur}]
from mcp.server import Server
from mcp.server.stdio import stdio_server
from mcp.types import Tool, TextContent

server = Server("local-supermemory")
store = MemoryStore()

@server.list_tools()
async def list_tools() -> list[Tool]:
    return [Tool(name="memory", inputSchema={...}), ...]

@server.call_tool()
async def call_tool(name: str, args: dict) -> list[TextContent]:
    if name == "memory":
        return [TextContent(type="text", text="...")]

async def run():
    async with stdio_server() as (r, w):
        await server.run(r, w, server.create_initialization_options())
\end{lstlisting}

\textbf{Kommunikation:} stdio (Subprocess), JSON-RPC, alle Handler async

%──────────────────────────────────────────────────────────────
\section{Erweiterung}
%──────────────────────────────────────────────────────────────

\subsection{Neues Tool}

\begin{lstlisting}[style=code,language=Python]
# 1. In list_tools() registrieren
Tool(name="export", description="Export als JSON", inputSchema={})

# 2. In call_tool() implementieren
elif name == "export":
    mems = store.get_all()
    return [TextContent(type="text", text=json.dumps(mems))]
\end{lstlisting}

\subsection{Custom Embedding}

\begin{lstlisting}[style=code,language=Python]
from chromadb.utils import embedding_functions
ef = embedding_functions.SentenceTransformerEmbeddingFunction(
    model_name="paraphrase-multilingual-MiniLM-L12-v2"
)
self.chroma.get_or_create_collection(name=name, embedding_function=ef)
\end{lstlisting}


%──────────────────────────────────────────────────────────────
\section{Troubleshooting}
%──────────────────────────────────────────────────────────────

\begin{description}
    \item[Server startet nicht] Python $\geq$3.10 prüfen, \texttt{pip install -e .}
    \item[Kein Profil] Ollama fehlt: \texttt{ollama pull qwen2.5:32b}
    \item[Leere Suche] Projekt-Name prüfen, Collection leer?
    \item[ChromaDB Fehler] Reset: \texttt{rm -rf \textasciitilde/.local-supermemory/chroma}
\end{description}

%──────────────────────────────────────────────────────────────
\section{Cloud vs. Local}
%──────────────────────────────────────────────────────────────

\begin{tabular}{lcc}
\toprule
\textbf{Feature} & \textbf{Cloud} & \textbf{Local} \\
\midrule
Speicher & Cloudflare & \texttt{\textasciitilde/.local-supermemory} \\
Datenschutz & Cloud & 100\% lokal \\
Offline & Nein & Ja \\
Embedding & Proprietär & MiniLM-L6-v2 \\
Profil-LLM & Cloud-API & Ollama \\
Kosten & Subscription & Kostenlos \\
\bottomrule
\end{tabular}

\vfill
\begin{center}
\rule{0.5\textwidth}{0.4pt}\\[1em]
\url{https://github.com/granaria/local-supermemory}
\end{center}

\end{document}
